\subsection{Datasets}

We evaluate on three publicly available industrial PdM datasets. Table~\ref{tab:datasets} summarizes their characteristics.

% Table is defined in paper/tables/table2_datasets.tex
\begin{table}[htbp]
\caption{Industrial PdM Dataset Summary}
\label{tab:datasets}
\centering
\begin{tabular}{lccccc}
\toprule
Dataset & Domain & Samples & Channels & Tasks & Failure Rate \\
\midrule
C-MAPSS & Turbofan & 100 & 21 & RUL/Forecast & 20\% \\
PHM Milling & CNC & 16k & 8 & Anomaly & 15\% \\
PU Bearings & Rotating & 32 & 4 & RUL/Anomaly & 32\% \\
Wind SCADA & Turbines & 6k & 52 & Forecast & 8\% \\
MIMII & Factory & 10k & 8 & Anomaly & 10\% \\
PRONOSTIA & Bearings & 17 & 3 & RUL & 100\% \\
\bottomrule
\end{tabular}
\end{table}


The \textbf{C-MAPSS} dataset~\cite{saxena2008cmapss} contains simulated run-to-failure trajectories of turbofan engines under four operating conditions (FD001--FD004) with 21 sensor channels. \textbf{Wind SCADA} captures operational parameters from wind turbines with 52 channels including power output, wind speed, and nacelle temperatures, enabling multivariate forecasting evaluation. \textbf{PHM Milling}~\cite{phm2010milling} contains high-speed CNC milling runs with force, vibration, and acoustic-emission measurements recorded for successive tool cuts. In our pipeline, each raw cut file is summarized into compact statistical features so the resulting cut-level sequence remains memory-safe on commodity Colab hardware.

\subsection{Models}

We evaluate three open-source TSFMs and one supervised baseline:
\begin{itemize}
    \item \textbf{MOMENT}~\cite{goswami2024moment}: 385M parameters, masked pretraining, multivariate support
    \item \textbf{Chronos}~\cite{ansari2024chronos}: Token-based approach inspired by LLMs, univariate, probabilistic
    \item \textbf{Lag-Llama}~\cite{lag_llama2023}: 7M parameters, autoregressive LLaMA-based forecasting, univariate, probabilistic
    \item \textbf{PatchTST}~\cite{nie2022patchtst}: Supervised patch-based transformer baseline (ICLR 2023)
\end{itemize}

Table~\ref{tab:capabilities} summarizes the key capabilities of each model. All three TSFMs support zero-shot inference in our benchmark. We retain an exploratory few-shot adapter interface in the codebase, but only zero-shot and cross-condition results are considered fully validated until stable fine-tuning implementations are available.

\subsection{Preprocessing Pipeline}

Our preprocessing pipeline addresses key challenges in industrial time series data:

\textbf{Chronological Splits:} We enforce strict chronological train/validation/test splits (70/15/15) with no shuffling to prevent temporal data leakage.

\textbf{Per-Sensor Normalization:} Z-score normalization per sensor, fitting statistics exclusively on training data.

\textbf{Missing Value Imputation:} Forward-fill followed by linear interpolation, with column-mean fallback for persistent gaps.

\textbf{Sequence Construction:} Sliding windows are constructed per trajectory with chronological split boundaries preserved. Wind SCADA uses lookback 512 and forecast horizon 96, while PHM Milling uses a shorter cut-level context (128 / 16) after streaming feature summarization.

\textbf{RUL Labeling:} Piecewise linear RUL labels capped at 125 cycles following the standard C-MAPSS protocol.

\subsection{Evaluation Protocol}

We evaluate models across two primary scenarios:

\textbf{Zero-Shot:} Direct inference without any task-specific adaptation, testing the transfer capability of TSFMs on unseen industrial data.

\textbf{Exploratory Few-Shot:} Adapter-based fine-tuning using 1\% of training data is retained as a future extension, but should only be reported in the final manuscript once the open-source wrappers perform real parameter updates.

We additionally report \textbf{cross-condition transfer} on C-MAPSS, training on FD001 and evaluating on FD002--FD004.

Metrics reported: MAE, RMSE, and MAPE for forecasting-style evaluations. Chronos is univariate; for multivariate inputs we average channels before inference and broadcast the forecast back to channel dimensions for a controlled comparison.
